As Large Language Models (LLMs) increasingly dominate automated text generation, the ability to hide stealthy signals within synthetic corpora becomes a significant security and forensic concern. While high-capacity steganography has been extensively studied, "diluted" steganography—where information is sparsely distributed based on rare linguistic triggers—remains a critical blind spot. We propose \ours, a framework to systematically evaluate the impact of signal dilution on LLM-based detection and extraction. We investigate trigger-word frequencies ranging from 7.8\% (dense) to 0.076\% (very sparse) in the \wikitext dataset. Our empirical evaluation using \gpt reveals that diluted steganography is effectively invisible to zero-shot detection, yielding a 0\% True Positive Rate across all dilution levels. Furthermore, while \gpt can extract bits from sparse signals with moderate accuracy ($\sim$71\%), it fails significantly as signal density increases due to the cognitive bottleneck of tracking patterns across long context windows. These findings suggest that signal dilution acts as a potent defense against automated detection, while sequence complexity serves as a natural barrier to reliable extraction.
